% model_extended.tex — Extended Mathematical Framework
% Capital Stack, Priority Waterfall, Preferred Stock, Conversion Optionality
%
% Extends the base model in plan.tex (BTC dynamics, OU premium, holding flywheel,
% ILE/TEE/PMRI/IBGR/IFRD) with Strategy's multi-layered capital structure.

\section{Extended Structural Model: Capital Stack and Preferred Equity}\label{sec:model_extended}

We extend the single-debt baseline of Section~2 to accommodate Strategy's
multi-layered capital structure, which as of early 2026 includes senior debt,
convertible bonds, multiple series of perpetual preferred stock (both
convertible and non-convertible), and common equity.  The extensions below
preserve all state variables and dynamics from the base model while enriching
the liability side and its implications for leverage, survival, and BTC
accretion.

%% ====================================================================
\subsection{Capital Stack and Priority Waterfall}\label{subsec:waterfall}

\begin{definition}[Capital Stack]\label{def:capital_stack}
At time $t$ the firm's capital structure comprises:
\begin{enumerate}
  \item \textbf{Senior debt} $D_s(t)$: term loans, secured notes, and
        operating liabilities with contractual priority.
  \item \textbf{Convertible debt} $D_c(t)$: unsecured convertible notes
        with face values, coupons, maturities, and equity conversion rights.
  \item \textbf{Non-convertible perpetual preferred stock} $P_{nc}(t)$:
        par value of outstanding non-convertible preferred (e.g.\ STRF at
        10\% cumulative dividend).
  \item \textbf{Convertible perpetual preferred stock} $P_c(t)$: par value
        of outstanding convertible preferred (e.g.\ STRK at 8\% cumulative
        dividend, convertible to 0.1 MSTR per share).
  \item \textbf{Common equity} $E(t)$: residual claim.
\end{enumerate}
\end{definition}

\noindent Total liabilities senior to common equity are
\begin{equation}\label{eq:total_senior}
  L_t = D_s(t) + D_c(t) + P_{nc}(t) + P_c(t).
\end{equation}

\begin{definition}[Liquidation Waterfall]\label{def:waterfall}
In a hypothetical liquidation at time $t$ with asset value
$A_t = H_t S_t$, claims are satisfied in strict priority:
\begin{equation}\label{eq:waterfall}
  \underbrace{D_s}_{\text{senior debt}}
  \;\succ\;
  \underbrace{D_c}_{\text{convertible debt}}
  \;\succ\;
  \underbrace{P_{nc}}_{\text{non-conv.\ pfd.}}
  \;\succ\;
  \underbrace{P_c}_{\text{conv.\ pfd.}}
  \;\succ\;
  \underbrace{E}_{\text{common}}.
\end{equation}
Each layer $\ell$ receives $\min(A_t^{(\ell)},\, C_\ell)$, where $C_\ell$ is
the face claim and $A_t^{(\ell)}$ is the residual asset value available after
all senior layers are paid:
\begin{align}
  A_t^{(1)} &= A_t, &
  R_s(t)     &= \min\bigl(A_t^{(1)},\, D_s(t)\bigr), \label{eq:wf_senior}\\
  A_t^{(2)} &= \bigl(A_t^{(1)} - D_s(t)\bigr)^+, &
  R_{dc}(t)  &= \min\bigl(A_t^{(2)},\, D_c(t)\bigr), \\
  A_t^{(3)} &= \bigl(A_t^{(2)} - D_c(t)\bigr)^+, &
  R_{nc}(t)  &= \min\bigl(A_t^{(3)},\, P_{nc}(t)\bigr), \\
  A_t^{(4)} &= \bigl(A_t^{(3)} - P_{nc}(t)\bigr)^+, &
  R_c(t)     &= \min\bigl(A_t^{(4)},\, P_c(t)\bigr), \\
  A_t^{(5)} &= \bigl(A_t^{(4)} - P_c(t)\bigr)^+, &
  E_t^{\text{res}} &= A_t^{(5)}. \label{eq:wf_equity}
\end{align}
\end{definition}

\noindent The base-model NAV generalizes to
\begin{equation}\label{eq:nav_extended}
  \mathrm{NAV}_t = \bigl(A_t - L_t\bigr)^+
  = \bigl(H_t S_t - D_s(t) - D_c(t) - P_{nc}(t) - P_c(t)\bigr)^+,
\end{equation}
and the log-premium identity becomes
\begin{equation}\label{eq:pi_extended}
  \pi_t = \log\!\left(\frac{E_t}{\mathrm{NAV}_t^{\mathrm{clip}}}\right),
  \qquad
  \mathrm{NAV}_t^{\mathrm{clip}} = \max\bigl(A_t - L_t,\, \varepsilon\bigr).
\end{equation}

\begin{remark}
For periods where some convertible debt is deep in-the-money, one may define
an \emph{effective liability}
$L_t^{\mathrm{eff}} = D_s(t) + D_c^{\mathrm{eff}}(t) + P_{nc}(t) + P_c(t)$,
where $D_c^{\mathrm{eff}}(t)$ discounts the equity-like component
(cf.\ Section~2.5 of the base model).  This keeps the waterfall
well-ordered while avoiding double-counting of the conversion value.
\end{remark}

%% ====================================================================
\subsection{Perpetual Preferred Stock Dynamics}\label{subsec:preferred}

Strategy has issued multiple series of perpetual preferred stock with fixed
cumulative dividends.  Unlike debt, preferred stock has no maturity; however,
failure to pay dividends triggers cumulative arrearages and potentially
accelerated conversion or governance consequences.

\begin{definition}[Preferred Dividend Obligations]\label{def:pfd_divs}
Let $\mathcal{P} = \{p_1,\ldots,p_M\}$ denote the $M$ outstanding preferred
series, each with par value $F_{p_i}$, number of shares $n_{p_i}$, and
contractual annual dividend rate $\delta_{p_i}$.  The aggregate annual
preferred dividend obligation is
\begin{equation}\label{eq:agg_pfd_div}
  \Phi_t = \sum_{i=1}^{M} \delta_{p_i}\, F_{p_i}\, n_{p_i}(t).
\end{equation}
For the current capital structure, the dominant terms are:
\begin{itemize}
  \item STRK: $\delta_{\mathrm{STRK}} = 0.08$, convertible.
  \item STRF: $\delta_{\mathrm{STRF}} = 0.10$, non-convertible.
\end{itemize}
\end{definition}

\begin{definition}[Dividend Coverage Ratio]\label{def:dcr}
The \emph{Dividend Coverage Ratio} at time $t$ measures the firm's capacity
to service preferred dividends from BTC-related cash flows:
\begin{equation}\label{eq:dcr}
  \mathrm{DCR}_t
  = \frac{\Delta A_t^+ + \mathrm{CF}_t^{\mathrm{ops}}}{\Phi_t\,\Delta t},
\end{equation}
where $\Delta A_t^+ = \max(\Delta A_t,\, 0)$ captures realized BTC
appreciation over the period $[t, t+\Delta t]$ and
$\mathrm{CF}_t^{\mathrm{ops}}$ is operating cash flow from the legacy
software business.  The firm is in \emph{dividend stress} when
$\mathrm{DCR}_t < 1$.
\end{definition}

\begin{proposition}[Extended Distress Threshold]\label{prop:distress}
Under the extended capital stack, the distress threshold (replacing the
base-model condition $A_t \le (1+\varepsilon)D_t$) incorporates both
insolvency and dividend-default risk:
\begin{equation}\label{eq:distress_ext}
  A_t \;<\; D_s(t) + D_c(t) + \mathrm{PV}_t\!\bigl(\text{preferred dividends}\bigr)
  + \varepsilon\bigl(P_{nc}(t) + P_c(t)\bigr),
\end{equation}
where the present value of preferred dividends over a horizon $T^*$
(e.g.\ 10 years) at discount rate $r$ is
\begin{equation}\label{eq:pv_pfd}
  \mathrm{PV}_t\!\bigl(\text{preferred dividends}\bigr)
  = \Phi_t \cdot \frac{1 - e^{-r T^*}}{r}.
\end{equation}
\end{proposition}

\begin{proof}
A firm that cannot meet both its hard liabilities (debt) and the present
value of its soft liabilities (cumulative preferred dividends) will either
(i) default on preferred dividends, triggering governance and market
consequences, or (ii) be forced into asset liquidation.  The buffer
$\varepsilon(P_{nc}+P_c)$ accounts for transaction costs and par-value
recovery expectations.
\end{proof}

\begin{definition}[Extended Survival Probability]\label{def:survival_ext}
Let the extended distress boundary be
\begin{equation}\label{eq:distress_boundary}
  \mathcal{B}_t = D_s(t) + D_c(t)
  + \mathrm{PV}_t\!\bigl(\text{preferred dividends}\bigr)
  + \varepsilon\bigl(P_{nc}(t) + P_c(t)\bigr).
\end{equation}
Define the first-passage time
\begin{equation}\label{eq:tau_ext}
  \tau^{\mathrm{ext}} = \inf\bigl\{u \in [0,T] : A_u \le \mathcal{B}_u\bigr\}.
\end{equation}
The extended survival probability is
\begin{equation}\label{eq:surv_ext}
  \mathbb{P}\!\left(\tau^{\mathrm{ext}} > T\right)
  \approx
  \frac{\#\{\text{simulated paths with no hit before } T\}}
       {\#\{\text{total paths}\}}.
\end{equation}
\end{definition}

\noindent The extended survival probability is strictly lower than the
base-model survival probability for any positive preferred dividend
obligation, reflecting the additional drag of perpetual fixed-income
commitments on BTC-leveraged assets.

%% ====================================================================
\subsection{Conversion Optionality}\label{subsec:conversion}

STRK preferred stock carries an embedded conversion right: each share may
be converted into $\phi = 0.1$ shares of MSTR common at the holder's
election.  This creates interplay between the preferred and equity layers.

\begin{definition}[Conversion Terms]\label{def:conversion}
Let $n_c(t)$ denote the number of outstanding convertible preferred shares
(STRK) at time $t$, with conversion ratio $\phi$ (shares of MSTR common
per preferred share).  The \emph{implied conversion strike price} is
\begin{equation}\label{eq:strike}
  K_c = \frac{F_c}{\phi},
\end{equation}
where $F_c$ is the liquidation preference (par value) per preferred share.
\end{definition}

\begin{proposition}[Conversion as Embedded American Option]\label{prop:conv_option}
Each STRK share embeds an American-style call option on MSTR common equity.
The holder converts when the equity value received exceeds the present value
of retaining the preferred position:
\begin{equation}\label{eq:conv_decision}
  \text{Convert at } t \quad\Longleftrightarrow\quad
  \phi \cdot P_t > F_c + \frac{\delta_c\, F_c}{r_c},
\end{equation}
where $r_c$ is the holder's required yield on the preferred, and the
right-hand side represents the par value plus the present value of the
perpetual dividend stream foregone upon conversion.  In a simplified
``intrinsic value'' framework:
\begin{equation}\label{eq:conv_intrinsic}
  \text{Convert when } P_t > K_c = \frac{F_c}{\phi}.
\end{equation}
\end{proposition}

\begin{remark}
In practice, because the preferred dividend provides ongoing income, rational
holders delay conversion beyond the intrinsic-value trigger.  The optimal
conversion boundary $P^*(t)$ solves a free-boundary problem analogous to the
American call with continuous dividends.  For simulation purposes, we use
the simplified threshold~\eqref{eq:conv_intrinsic} with an optional
``conversion premium'' multiplier $\psi \ge 1$:
\begin{equation}\label{eq:conv_trigger}
  P_t > \psi \cdot K_c.
\end{equation}
\end{remark}

\begin{definition}[Post-Conversion Dynamics]\label{def:post_conv}
When a fraction $\xi_t \in [0,1]$ of convertible preferred shares convert at
time $t$:
\begin{align}
  N_t^{\mathrm{sh}} &\;\longrightarrow\;
    N_t^{\mathrm{sh}} + \xi_t\, n_c(t)\, \phi,
    \label{eq:dilution}\\
  n_c(t) &\;\longrightarrow\; (1-\xi_t)\, n_c(t),
    \label{eq:pfd_reduction}\\
  \Phi_t &\;\longrightarrow\;
    \Phi_t - \xi_t\, \delta_c\, F_c\, n_c(t).
    \label{eq:div_reduction}
\end{align}
Conversion simultaneously \emph{dilutes} common shares
\eqref{eq:dilution}, \emph{reduces} the preferred overhang
\eqref{eq:pfd_reduction}, and \emph{lowers} dividend obligations
\eqref{eq:div_reduction}.
\end{definition}

\begin{definition}[Conversion-Adjusted NAV]\label{def:ca_nav}
After accounting for potential conversion of all in-the-money convertible
preferred shares, the \emph{Conversion-Adjusted NAV} is
\begin{equation}\label{eq:ca_nav}
  \mathrm{NAV}_t^{\mathrm{CA}}
  = \frac{A_t - D_s(t) - D_c(t) - P_{nc}(t)}
         {N_t^{\mathrm{sh}} + n_c^{\mathrm{ITM}}(t)\,\phi},
\end{equation}
where $n_c^{\mathrm{ITM}}(t)$ is the number of convertible preferred shares
for which $\phi\, P_t > F_c$ (i.e.\ in-the-money).
\end{definition}

\begin{definition}[Conversion-Adjusted BTC per Share]\label{def:ca_bps}
The \emph{Conversion-Adjusted BTC per Share} accounts for the maximum
dilution from all in-the-money conversions:
\begin{equation}\label{eq:ca_bps}
  B_t^{\mathrm{CA}}
  = \frac{H_t}{N_t^{\mathrm{sh}} + n_c^{\mathrm{ITM}}(t)\,\phi}.
\end{equation}
This provides a conservative (fully diluted) measure of BTC ownership per
common share.
\end{definition}

%% ====================================================================
\subsection{Multi-Channel Financing Flywheel}\label{subsec:flywheel}

The base model captures BTC accumulation via the holding dynamics
$dH_t = \alpha\,\pi_t^+\,dt + dH_t^{\mathrm{jump}}$.
We extend this to disaggregate the financing channels.

\begin{definition}[Multi-Channel Holding Dynamics]\label{def:multichannel}
BTC accumulation is decomposed by financing source:
\begin{equation}\label{eq:dH_multi}
  dH_t = dH_t^{\mathrm{ATM}} + dH_t^{\mathrm{conv}} + dH_t^{\mathrm{pfd}}
         + dH_t^{\mathrm{ops}},
\end{equation}
where:
\begin{enumerate}
  \item $dH_t^{\mathrm{ATM}}$: BTC acquired from at-the-market common equity
        issuance.  Proceeds $= (\text{shares issued}) \times P_t$; BTC
        purchased $= \text{proceeds} / S_t$.
  \item $dH_t^{\mathrm{conv}}$: BTC acquired from convertible debt issuance.
        Proceeds $= D_c^{\mathrm{new}}(t)$; BTC purchased
        $= D_c^{\mathrm{new}}(t) / S_t$.
  \item $dH_t^{\mathrm{pfd}}$: BTC acquired from preferred stock issuance.
        Proceeds $= P^{\mathrm{new}}(t)$; BTC purchased
        $= P^{\mathrm{new}}(t) / S_t$.
  \item $dH_t^{\mathrm{ops}}$: BTC acquired from operating cash flows
        (typically small).
\end{enumerate}
\end{definition}

Each channel entails different cost of capital, dilution characteristics,
and market conditions under which it is accessible.

\begin{proposition}[Channel Selection]\label{prop:channel_select}
The optimal financing channel depends on the premium $\pi_t$ and market
conditions.  We model channel intensities as functions of state:
\begin{align}
  \lambda_t^{\mathrm{ATM}} &= \lambda_0^{\mathrm{ATM}}\,
    \mathbf{1}\{\pi_t > \pi^{\mathrm{ATM}}_{\min}\}\,
    g_{\mathrm{ATM}}(\pi_t),
    \label{eq:lambda_atm}\\
  \lambda_t^{\mathrm{conv}} &= \lambda_0^{\mathrm{conv}}\,
    \mathbf{1}\{\sigma_S(t) > \bar\sigma\}\,
    g_{\mathrm{conv}}(\pi_t, r_t),
    \label{eq:lambda_conv}\\
  \lambda_t^{\mathrm{pfd}} &= \lambda_0^{\mathrm{pfd}}\,
    g_{\mathrm{pfd}}(\pi_t, \mathrm{DCR}_t),
    \label{eq:lambda_pfd}
\end{align}
where $g_{\cdot}(\cdot)$ are non-negative scaling functions and the
indicators encode minimum thresholds: ATM issuance requires a positive
premium, convertible issuance benefits from high BTC volatility (making
the embedded option more valuable), and preferred issuance is viable when
dividend coverage is adequate.
\end{proposition}

\begin{definition}[Weighted Average Cost of BTC Acquisition (WACBA)]%
\label{def:wacba}
The \emph{Weighted Average Cost of BTC Acquisition} over the period
$[0, T]$ accounts for the blended cost of capital across all financing
channels:
\begin{equation}\label{eq:wacba}
  \mathrm{WACBA}_{[0,T]}
  = \frac{\sum_{j \in \mathcal{C}}
          c_j \displaystyle\int_0^T S_u\,dH_u^{(j)}}
         {\displaystyle\int_0^T dH_u},
\end{equation}
where $\mathcal{C} = \{\mathrm{ATM},\,\mathrm{conv},\,\mathrm{pfd},\,
\mathrm{ops}\}$ indexes the channels and $c_j$ is the channel-specific
cost multiplier that captures dilution and ongoing obligations:
\begin{align}
  c_{\mathrm{ATM}} &= 1 + \frac{\Delta N^{\mathrm{sh}}_{\mathrm{ATM}}}
                               {N_t^{\mathrm{sh}}}\cdot
                         \frac{H_t\,S_t}{\text{ATM proceeds}},
                         \label{eq:c_atm}\\
  c_{\mathrm{conv}} &= 1 + \frac{\text{coupon PV} + \text{conversion dilution PV}}
                                {\text{convertible proceeds}},
                                \label{eq:c_conv}\\
  c_{\mathrm{pfd}} &= 1 + \frac{\mathrm{PV}(\text{preferred dividends})}
                               {\text{preferred proceeds}},
                               \label{eq:c_pfd}\\
  c_{\mathrm{ops}} &= 1. \label{eq:c_ops}
\end{align}
\end{definition}

\begin{remark}
The WACBA generalizes the ``average BTC purchase price'' reported by
Strategy.  Channels with lower $c_j$ are more efficient per BTC acquired.
ATM issuance at high premiums can have $c_{\mathrm{ATM}} < 1$ in
BTC-per-share terms (accretive issuance), while preferred issuance carries
a perpetual drag via $c_{\mathrm{pfd}}$.
\end{remark}

%% ====================================================================
\subsection{Dilution-Adjusted Dynamics}\label{subsec:dilution_adj}

Each financing channel affects the share count differently.  We generalize
the share-count process from the base model:
\begin{equation}\label{eq:dN_extended}
  dN_t^{\mathrm{sh}}
  = dN_t^{\mathrm{ATM}} + dN_t^{\mathrm{conv\_debt}} + dN_t^{\mathrm{conv\_pfd}},
\end{equation}
where:
\begin{itemize}
  \item $dN_t^{\mathrm{ATM}} = (\text{ATM proceeds at }t) / P_t$: shares
        issued via at-the-market program.
  \item $dN_t^{\mathrm{conv\_debt}}$: shares created upon conversion of
        convertible bonds.
  \item $dN_t^{\mathrm{conv\_pfd}} = \xi_t\, n_c(t)\, \phi$: shares created
        upon conversion of convertible preferred (STRK).
\end{itemize}

\begin{proposition}[BTC-per-Share Accretion Condition]\label{prop:accretion}
ATM issuance at time $t$ is BTC-per-share \emph{accretive} if and only if
the implied BTC yield per dollar raised exceeds the current BTC per share
loading:
\begin{equation}\label{eq:accretive}
  \frac{1}{S_t} > \frac{B_t}{P_t}
  \quad\Longleftrightarrow\quad
  P_t > B_t \cdot S_t
  \quad\Longleftrightarrow\quad
  \frac{P_t\, N_t^{\mathrm{sh}}}{H_t\, S_t} > 1
  \quad\Longleftrightarrow\quad
  \pi_t > 0.
\end{equation}
That is, ATM issuance is accretive precisely when the firm trades at a
premium to NAV.
\end{proposition}

\begin{proof}
ATM issuance of $\Delta N$ shares at price $P_t$ raises $P_t\,\Delta N$
dollars, purchasing $P_t\,\Delta N / S_t$ BTC.  Post-issuance BTC per share
is
\[
  B_t^{\mathrm{post}}
  = \frac{H_t + P_t\,\Delta N / S_t}{N_t^{\mathrm{sh}} + \Delta N}.
\]
$B_t^{\mathrm{post}} > B_t$ iff $P_t / S_t > B_t = H_t / N_t^{\mathrm{sh}}$,
which is equivalent to $E_t > H_t S_t - D_t$, i.e.\ $\pi_t > 0$ when
$\mathrm{NAV}_t > 0$.
\end{proof}

%% ====================================================================
\subsection{Extended Indicators}\label{subsec:new_indicators}

We define five new indicators that complement the base-model suite
(IBGR, ILE, TEE, PMRI, IFRD).

\subsubsection{Capital Stack Leverage Ratio (CSLR)}

\begin{definition}[CSLR]\label{def:cslr}
The \emph{Capital Stack Leverage Ratio} measures total claims senior to
common equity relative to equity market value:
\begin{equation}\label{eq:cslr}
  \mathrm{CSLR}_t
  = \frac{D_s(t) + D_c(t) + P_{nc}(t) + P_c(t)}{E_t}
  = \frac{L_t}{E_t}.
\end{equation}
\end{definition}

\noindent Unlike the base-model ILE, which captures \emph{elasticity},
the CSLR is a static leverage ratio that directly measures the weight
of the liability stack.  As preferred stock issuance grows, CSLR rises
even if debt remains constant, capturing the layering of soft liabilities.

\subsubsection{Preferred Coverage Ratio (PCR)}

\begin{definition}[PCR]\label{def:pcr}
The \emph{Preferred Coverage Ratio} measures the firm's capacity to
meet preferred dividend obligations:
\begin{equation}\label{eq:pcr}
  \mathrm{PCR}_t
  = \frac{\mathbb{E}_t^{\mathbb{P}}\!\bigl[\Delta A_{t,T}^+\bigr]
          + \mathrm{CF}_t^{\mathrm{ops}}\cdot T}
         {\Phi_t \cdot T},
\end{equation}
where $\Delta A_{t,T}^+ = \max(A_T - A_t,\, 0)$ is expected BTC
appreciation over horizon $T$, and $\Phi_t$ is the annual preferred
dividend obligation from~\eqref{eq:agg_pfd_div}.
\end{definition}

\noindent $\mathrm{PCR}_t > 1$ indicates that expected asset appreciation
covers preferred dividends; $\mathrm{PCR}_t < 1$ signals potential
dividend stress and increased probability of governance consequences
(e.g.\ board representation rights for preferred holders).

\subsubsection{Dilution-Adjusted IBGR}

\begin{definition}[Dilution-Adjusted IBGR]\label{def:da_ibgr}
The \emph{Dilution-Adjusted IBGR} measures BTC growth per share net of
all dilution sources, including potential preferred stock conversions:
\begin{equation}\label{eq:da_ibgr}
  \mathrm{IBGR}_t^{\mathrm{DA}}
  = \frac{1}{B_t^{\mathrm{CA}}}
    \mathbb{E}_t^{\mathbb{P}}\!\left[
      \frac{B_{t+\Delta t}^{\mathrm{CA}} - B_t^{\mathrm{CA}}}{\Delta t}
    \right],
\end{equation}
where $B_t^{\mathrm{CA}}$ is the Conversion-Adjusted BTC per Share
from~\eqref{eq:ca_bps}.  This differs from the base-model per-share IBGR
by using the fully diluted share count.
\end{definition}

\noindent In practice, $\mathrm{IBGR}_t^{\mathrm{DA}}$ is estimated via
Monte Carlo by simulating joint paths of $(S_t, H_t, N_t^{\mathrm{sh}},
n_c(t))$ under $\mathbb{P}$ and computing sample means of the ratio.

\subsubsection{Effective Leverage Elasticity (ELE)}

\begin{definition}[ELE]\label{def:ele}
The \emph{Effective Leverage Elasticity} generalizes ILE to account for
the preferred stock buffer layer between debt and common equity:
\begin{equation}\label{eq:ele}
  \mathrm{ELE}_t
  = \frac{\partial \log E_t}{\partial \log S_t}\bigg|_{\pi_t\text{ fixed}}
  = \frac{H_t\, S_t}
         {H_t\, S_t - D_s(t) - D_c(t) - P_{nc}(t) - P_c(t)}
  = \frac{A_t}{A_t - L_t},
\end{equation}
whenever $A_t > L_t$.
\end{definition}

\begin{proposition}[Ordering of Elasticities]\label{prop:ele_ordering}
When preferred stock outstanding is positive,
\begin{equation}\label{eq:ele_vs_ile}
  \mathrm{ELE}_t \;\ge\; \mathrm{ILE}_t
  = \frac{A_t}{A_t - D_s(t) - D_c(t)},
\end{equation}
with equality iff $P_{nc}(t) + P_c(t) = 0$.  The incremental leverage
from the preferred layer is
\begin{equation}\label{eq:ele_increment}
  \Delta\beta_t^{\mathrm{pfd}}
  = \mathrm{ELE}_t - \mathrm{ILE}_t
  = \frac{A_t\bigl(P_{nc}(t) + P_c(t)\bigr)}
         {\bigl(A_t - D_t\bigr)\bigl(A_t - L_t\bigr)},
\end{equation}
where $D_t = D_s(t) + D_c(t)$.
\end{proposition}

\begin{proof}
Both $\mathrm{ELE}_t$ and $\mathrm{ILE}_t$ have the form $A_t/(A_t - X)$
with $X_{\mathrm{ELE}} = L_t \ge X_{\mathrm{ILE}} = D_t$.  The function
$f(X) = A_t/(A_t-X)$ is increasing in $X$ for $X < A_t$, yielding the
inequality.  The difference follows by algebra.
\end{proof}

\noindent The extended Total Equity Elasticity becomes
\begin{equation}\label{eq:tee_ext}
  \beta_{\mathrm{TOT}}^{\mathrm{ext}}(t)
  = \gamma_{\pi S} + \mathrm{ELE}_t,
\end{equation}
replacing $\mathrm{ILE}_t$ with $\mathrm{ELE}_t$ in the base formula.

\subsubsection{Capital Efficiency Ratio (CER)}

\begin{definition}[CER]\label{def:cer}
The \emph{Capital Efficiency Ratio} measures the BTC-per-share accretion
generated per unit of capital stack leverage:
\begin{equation}\label{eq:cer}
  \mathrm{CER}_t
  = \frac{\mathrm{IBGR}_t^{\mathrm{DA}}}{\mathrm{CSLR}_t}.
\end{equation}
\end{definition}

\noindent High CER indicates that additional leverage is being deployed
productively (generating BTC accretion per share), while low CER signals
that the capital stack is growing faster than per-share BTC accumulation
--- a sign of diminishing returns to the flywheel strategy.

%% ====================================================================
\subsection{Structural Relationships and Summary}\label{subsec:summary}

\begin{theorem}[Flywheel Decomposition]\label{thm:flywheel_decomp}
Under the extended model, the instantaneous change in common equity value
decomposes as:
\begin{equation}\label{eq:equity_decomp}
  \frac{dE_t}{E_t}
  = \underbrace{\mathrm{ELE}_t \cdot \frac{dS_t}{S_t}}_{\text{leveraged BTC return}}
  + \underbrace{d\pi_t}_{\text{premium change}}
  - \underbrace{\frac{\Phi_t\,dt}{A_t - L_t}}_{\text{preferred dividend drag}}
  + \underbrace{\frac{S_t\,dH_t}{A_t - L_t}}_{\text{BTC accretion lift}}
  + O(dt).
\end{equation}
\end{theorem}

\begin{proof}
From $E_t = e^{\pi_t}(A_t - L_t)^+$ with $A_t = H_t S_t$, apply It\^{o}'s
lemma to $\log E_t = \pi_t + \log(A_t - L_t)$ and note that preferred
dividends reduce $A_t$ at rate $\Phi_t$ (dividends are paid from assets),
while BTC acquisitions increase $A_t$ at rate $S_t\,dH_t$.  Collecting
terms yields the decomposition.
\end{proof}

\noindent This decomposition makes explicit the tension at the heart of
the Strategy model:
\begin{itemize}
  \item The BTC accretion lift is powered by the financing flywheel
        (Section~\ref{subsec:flywheel}).
  \item The preferred dividend drag is the cost of funding via preferred
        equity, which grows monotonically with issuance.
  \item The premium change can amplify or dampen returns depending on
        market sentiment.
  \item Effective leverage $\mathrm{ELE}_t$ increases with both debt and
        preferred issuance, amplifying BTC returns in both directions.
\end{itemize}

Table~\ref{tab:indicators} summarizes the complete indicator suite.

\begin{table}[htbp]
\centering
\caption{Summary of structural indicators: base model and extensions.}
\label{tab:indicators}
\small
\begin{tabular}{@{}llp{6.5cm}@{}}
\toprule
\textbf{Indicator} & \textbf{Type} & \textbf{Description} \\
\midrule
ILE & Base & Structural leverage elasticity (debt only) \\
TEE & Base & Total equity elasticity incl.\ premium response \\
PMRI & Base & Premium z-score under OU stationarity \\
IBGR & Base & Expected BTC growth rate (total and per-share) \\
IFRD & Base & Funding requirement distribution \\
\midrule
CSLR & Extended & Total claims / equity (full capital stack leverage) \\
PCR & Extended & Expected asset appreciation / preferred dividends \\
$\text{IBGR}^{\text{DA}}$ & Extended & Dilution-adjusted BTC growth per share \\
ELE & Extended & Leverage elasticity incl.\ preferred buffer \\
CER & Extended & BTC accretion per unit leverage \\
\bottomrule
\end{tabular}
\end{table}
